\documentclass[11pt,a4paper]{article}
\usepackage[utf8]{inputenc}
\usepackage{amsmath, amssymb, amsthm}
\usepackage{booktabs}
\usepackage{geometry}
\usepackage{graphicx}
\usepackage{hyperref}

\geometry{margin=1in}

\title{\textbf{Lecture Notes: The Crisis of Classical Electromagnetism \\ The Poynting Vector vs. The Photoelectric Effect}}
\author{\textbf{Advanced Physics Seminar}}
\date{\today}

\begin{document}
	
	\maketitle
	
	\section{Introduction}
	In classical electromagnetism, light is modeled as a continuous electromagnetic wave where energy is distributed uniformly across an advancing wavefront. Conversely, the quantum mechanical description treats light as localized packets of energy called photons. This lecture explores the stark contradictions between these two frameworks, focusing on the mathematical breakdown of the classical \textbf{Poynting vector} when applied to the \textbf{photoelectric effect}, and its subsequent quantum resolution.
	
	---
	
	\section{The Classical Framework: The Poynting Vector}
	According to classical electromagnetic theory, the rate of energy transfer per unit area (energy flux density) by an electromagnetic wave is quantified by the \textbf{Poynting vector} ($\vec{S}$):
	\begin{equation}
		\vec{S} = \frac{1}{\mu_0} (\vec{E} \times \vec{B})
	\end{equation}
	Where:
	\begin{itemize}
		\item $\vec{E}$ and $\vec{B}$ are the oscillating electric and magnetic field vectors.
		\item $\mu_0$ is the permeability of free space.
	\end{itemize}
	
	For a monochromatic plane wave propagating in the $z$-direction, the fields oscillate as:
	\begin{align}
		E(t) &= E_0 \cos(\omega t - kz) \\
		B(t) &= B_0 \cos(\omega t - kz)
	\end{align}
	Substituting these into the instantaneous Poynting vector yields:
	\begin{equation}
		S(t) = c \epsilon_0 E_0^2 \cos^2(\omega t - kz) = \frac{1}{2} c \epsilon_0 E_0^2 \Big[ 1 + \cos(2\omega t - 2kz) \Big]
	\end{equation}
	
	\subsection{The Time-Averaged Intensity}
	Taking the time average over a full wave period ($\langle \dots \rangle$) eliminates the oscillating frequency component entirely:
	\begin{equation}
		\langle S \rangle = \frac{1}{2} c \epsilon_0 E_0^2
	\end{equation}
	In this classical framework, \textbf{frequency ($\omega$) and electric field amplitude ($E_0$) are completely independent variables}. You can modulate the peak force ($E_0$) to be massive while maintaining an incredibly low frequency.
	
	---
	
	\section{The Three Failures of Classical Theory}
	When applied to the photoelectric effect—where light ejects electrons from a metal surface—the classical model fails dramatically in three ways:
	
	\begin{table}[h!]
		\centering
		\caption{Classical Wave Theory vs. Quantum Reality}
		\begin{tabular}{p{7cm}p{8cm}}
			\toprule
			\textbf{Classical Wave Expectation ($\langle S \rangle$)} & \textbf{Experimental Quantum Reality} \\
			\midrule
			\textbf{1. The Time Delay Paradox:} Energy is smeared thin. Localized electrons must act like sponges, accumulating energy over hours or days before escaping. & \textbf{Instantaneous Emission:} Electrons are ejected almost immediately ($< 10^{-9}$ s), even under extremely faint light conditions. \\
			\addlinespace
			\textbf{2. The Intensity Problem:} Higher intensity means a larger amplitude ($E_0$). Taller wave crests apply larger physical forces ($\vec{F}=q\vec{E}$), predicting faster electrons. & \textbf{Kinetic Energy Independent of Intensity:} Increasing brightness only yields \textit{more} electrons. The maximum kinetic energy remains completely unaltered. \\
			\addlinespace
			\textbf{3. The Frequency Blind Spot:} If the wave is intense enough (huge $E_0^2$), it should deliver enough energy to eject electrons at \textit{any} frequency. & \textbf{Threshold Frequency ($\nu_0$):} Below a strict minimum frequency, zero electrons are emitted, regardless of how intense or long the light shines. \\
			\bottomrule
		\end{tabular}
	\end{table}
	
	\subsection{The Energy Peak Paradox}
	Even if one analyzes the \textit{instantaneous} Poynting vector $S(t)$, the maximum peak of the flux is $c \epsilon_0 E_0^2$. For high-intensity red light, this peak force is massive. Classical mechanics dictates that this violent jerk of electric field force should violently rip the electron out during the cycle's peak. 
	
	However, experiments show this never happens. Because the frequency is wrong, the electron simply rides the wave up and down symmetrically. It absorbs a push in the first half-cycle and receives an equal, opposite push in the second half-cycle, resulting in a net energy transfer of zero.
	
	---
	
	\section{Case Study: The Time-Delay Mathematical Proof}
	To illustrate the severity of the classical breakdown, let us calculate the classical waiting time for an electron on a \textbf{sodium plate}.
	
	\subsection{Given Parameters}
	\begin{itemize}
		\item Work Function of Sodium ($\Phi$): $2.28 \text{ eV} = 2.28 \times 1.6 \times 10^{-19} \text{ J} = 3.65 \times 10^{-19} \text{ J}$.
		\item Target Electron Cross-Sectional Area ($A$): Modeled as an atomic disk where $r \approx 0.1 \text{ nm}$.
		\begin{equation}
			A = \pi r^2 = \pi \times (1 \times 10^{-10} \text{ m})^2 \approx 3.14 \times 10^{-20} \text{ m}^2
		\end{equation}
		\item Light Intensity ($\langle S \rangle$): A standard laboratory value of $1.00 \text{ W/m}^2$.
	\end{itemize}
	
	\subsection{Calculation}
	Assuming the Poynting vector delivers energy continuously and uniformly across the surface, the power absorbed by a single electron target is:
	\begin{equation}
		P_{\text{absorbed}} = \langle S \rangle \times A = 1.00 \text{ W/m}^2 \times 3.14 \times 10^{-20} \text{ m}^2 = 3.14 \times 10^{-20} \text{ J/s}
	\end{equation}
	The time required for the electron to collect enough continuous energy to match the work function $\Phi$ is:
	\begin{equation}
		\Delta t = \frac{\Phi}{P_{\text{absorbed}}} = \frac{3.65 \times 10^{-19} \text{ J}}{3.14 \times 10^{-20} \text{ J/s}} \approx \mathbf{11.62 \text{ seconds}}
	\end{equation}
	\textit{Note: If the intensity is scaled down to a faint $0.01 \text{ W/m}^2$, the classical waiting time shoots up to nearly 20 minutes. Yet, experimentally, emission remains instantaneous.}
	
	---
	
	\section{The Quantum Resolution}
	Albert Einstein resolved the crisis in 1905 by fundamentally reinterpreting light's core parameters into a quantized framework, changing how the concepts enter the Poynting vector:
	
	\begin{enumerate}
		\item \textbf{Frequency ($\nu$) = The Payload:} Frequency determines the precise, indivisible energy packet of a single photon: $E_{\text{photon}} = h\nu$.
		\item \textbf{Intensity ($\langle S \rangle$) = The Photon Count:} Turning up the brightness does not make the wave peaks taller. It increases the \textit{photon flux density} ($N$, the number of photons crossing a unit area per second).
	\end{enumerate}
	
	Thus, the quantum equivalence of the average Poynting vector magnitude maps as:
	\begin{equation}
		\langle S \rangle = \frac{1}{2} c \epsilon_0 E_0^2 \equiv N \cdot (h\nu)
	\end{equation}
	Where the square of the field amplitude ($E_0^2$) acts as a \textbf{probability density operator} for finding photons, and the frequency governs the energetic punch. Because electrons and photons interact on a strict 1-to-1 collision basis, an electron cannot store up partial wave fractions. It either takes the entire payload or nothing.
	
	Einstein's conservation of energy formula governs this system perfectly:
	\begin{equation}
		K_{\text{max}} = h\nu - \Phi
	\end{equation}
	
	---
	
	\section{Experimental Verification: The Stopping Potential Apparatus}
	To prove Einstein's quantum model over the classical Poynting vector, physicists utilized a photoelectric tube experiment. 
	
	An evacuated quartz vacuum tube contains an Emitter (Target Metal) and a Collector plate. Monochromatic light knocks electrons free, generating a \textbf{photocurrent} measured by an ammeter.
	
	\subsection{The Electrical Brake Mechanism}
	To measure the exact maximum kinetic energy ($K_{\text{max}}$) of the escaping electrons, the voltage polarity is reversed. The Collector is made negative relative to the Emitter. This establishes a retarding electric field that repels incoming electrons.
	
	The negative voltage is increased until the ammeter drops to exactly zero. This critical voltage is the \textbf{Stopping Potential ($V_0$)}. By conservation of energy, the work done by the electric field to stop the fastest electron equals its initial kinetic energy:
	\begin{equation}
		K_{\text{max}} = e V_0
	\end{equation}
	Substituting this into Einstein's equation gives the linear relationship:
	\begin{equation}
		e V_0 = h\nu - \Phi \implies V_0 = \left(\frac{h}{e}\right)\nu - \frac{\Phi}{e}
	\end{equation}
	
	\subsection{Experimental Conclusions}
	\begin{itemize}
		\item \textbf{Varying Intensity:} Changing the intensity (modulating $\langle S \rangle$ via $E_0^2$) shifts the total current magnitude but leaves $V_0$ completely unchanged. 
		\item \textbf{Varying Frequency:} Modulating the frequency ($\nu$) shifts the stopping potential $V_0$ linearly, with the slope of the line yielding the fundamental constant ratio $\frac{h}{e}$. This finalized the proof that light transfers energy via discrete, frequency-locked packets.
	\end{itemize}
	
\end{document}
